\documentclass[
  man,
  floatsintext,
  longtable,
  nolmodern,
  notxfonts,
  notimes,
  12pt,
  colorlinks=true,linkcolor=blue,citecolor=blue,urlcolor=blue]{apa7}

\usepackage{amsmath}
\usepackage{amssymb}
\usepackage{setspace}


\usepackage[bidi=default]{babel}
\babelprovide[main,import]{english}


% get rid of language-specific shorthands (see #6817):
\let\LanguageShortHands\languageshorthands
\def\languageshorthands#1{}

\RequirePackage{longtable}
\RequirePackage{threeparttablex}

\usepackage{framed}


\makeatletter
\renewcommand{\paragraph}{\@startsection{paragraph}{4}{\parindent}%
	{0\baselineskip \@plus 0.2ex \@minus 0.2ex}%
	{-.5em}%
	{\normalfont\normalsize\bfseries\typesectitle}}

\renewcommand{\subparagraph}[1]{\@startsection{subparagraph}{5}{0.5em}%
	{0\baselineskip \@plus 0.2ex \@minus 0.2ex}%
	{-\z@\relax}%
	{\normalfont\normalsize\bfseries\itshape\hspace{\parindent}{#1}\textit{\addperi}}{\relax}}
\makeatother




% longtable is already loaded via doc-class.tex (apa7 class option) and
% apatemplate.tex (\RequirePackage{longtable}). Keep booktabs here so raw
% LaTeX tables using \toprule/\midrule still render even without Pandoc table
% nodes.
\usepackage{booktabs, multirow, multicol, colortbl, hhline, caption, array, float, xpatch}
\usepackage{subcaption}


\renewcommand\thesubfigure{\Alph{subfigure}}
\setcounter{topnumber}{2}
\setcounter{bottomnumber}{2}
\setcounter{totalnumber}{4}
\renewcommand{\topfraction}{0.85}
\renewcommand{\bottomfraction}{0.85}
\renewcommand{\textfraction}{0.15}
\renewcommand{\floatpagefraction}{0.7}

\usepackage{tcolorbox}
\tcbuselibrary{listings,theorems, breakable, skins}
\usepackage{fontawesome5}

\definecolor{quarto-callout-color}{HTML}{909090}
\definecolor{quarto-callout-note-color}{HTML}{0758E5}
\definecolor{quarto-callout-important-color}{HTML}{CC1914}
\definecolor{quarto-callout-warning-color}{HTML}{EB9113}
\definecolor{quarto-callout-tip-color}{HTML}{00A047}
\definecolor{quarto-callout-caution-color}{HTML}{FC5300}
\definecolor{quarto-callout-color-frame}{HTML}{ACACAC}
\definecolor{quarto-callout-note-color-frame}{HTML}{4582EC}
\definecolor{quarto-callout-important-color-frame}{HTML}{D9534F}
\definecolor{quarto-callout-warning-color-frame}{HTML}{F0AD4E}
\definecolor{quarto-callout-tip-color-frame}{HTML}{02B875}
\definecolor{quarto-callout-caution-color-frame}{HTML}{FD7E14}

%\newlength\Oldarrayrulewidth
%\newlength\Oldtabcolsep


\usepackage{hyperref}
\usepackage[protrusion=true]{microtype}


\providecommand{\tightlist}{%
  \setlength{\itemsep}{0pt}\setlength{\parskip}{0pt}}
\usepackage{longtable,booktabs,array}
\newcounter{none} % for unnumbered tables
\usepackage{calc} % for calculating minipage widths
% Correct order of tables after \paragraph or \subparagraph
\usepackage{etoolbox}
\makeatletter
\patchcmd\longtable{\par}{\if@noskipsec\mbox{}\fi\par}{}{}
\makeatother
% Allow footnotes in longtable head/foot
\IfFileExists{footnotehyper.sty}{\usepackage{footnotehyper}}{\usepackage{footnote}}
\makesavenoteenv{longtable}

\usepackage{graphicx}
\makeatletter
\newsavebox\pandoc@box
\newcommand*\pandocbounded[1]{% scales image to fit in text height/width
  \sbox\pandoc@box{#1}%
  \Gscale@div\@tempa{\textheight}{\dimexpr\ht\pandoc@box+\dp\pandoc@box\relax}%
  \Gscale@div\@tempb{\linewidth}{\wd\pandoc@box}%
  \ifdim\@tempb\p@<\@tempa\p@\let\@tempa\@tempb\fi% select the smaller of both
  \ifdim\@tempa\p@<\p@\scalebox{\@tempa}{\usebox\pandoc@box}%
  \else\usebox{\pandoc@box}%
  \fi%
}
% Set default figure placement to htbp
\def\fps@figure{htbp}
\makeatother


% definitions for citeproc citations
\NewDocumentCommand\citeproctext{}{}
\NewDocumentCommand\citeproc{mm}{%
  \begingroup\def\citeproctext{#2}\cite{#1}\endgroup}
\makeatletter
 % allow citations to break across lines
 \let\@cite@ofmt\@firstofone
 % avoid brackets around text for \cite:
 \def\@biblabel#1{}
 \def\@cite#1#2{{#1\if@tempswa , #2\fi}}
\makeatother
\newlength{\cslhangindent}
\setlength{\cslhangindent}{1.5em}
\newlength{\csllabelwidth}
\setlength{\csllabelwidth}{3em}
\newenvironment{CSLReferences}[2] % #1 hanging-indent, #2 entry-spacing
 {\begin{list}{}{%
  \setlength{\itemindent}{0pt}
  \setlength{\leftmargin}{0pt}
  \setlength{\parsep}{0pt}
  % turn on hanging indent if param 1 is 1
  \ifodd #1
   \setlength{\leftmargin}{\cslhangindent}
   \setlength{\itemindent}{-1\cslhangindent}
  \fi
  % set entry spacing
  \setlength{\itemsep}{#2\baselineskip}}}
 {\end{list}}
\usepackage{calc}
\newcommand{\CSLBlock}[1]{\hfill\break\parbox[t]{\linewidth}{\strut\ignorespaces#1\strut}}
\newcommand{\CSLLeftMargin}[1]{\parbox[t]{\csllabelwidth}{\strut#1\strut}}
\newcommand{\CSLRightInline}[1]{\parbox[t]{\linewidth - \csllabelwidth}{\strut#1\strut}}
\newcommand{\CSLIndent}[1]{\hspace{\cslhangindent}#1}





\usepackage{newtx}

\defaultfontfeatures{Scale=MatchLowercase}
\defaultfontfeatures[\rmfamily]{Ligatures=TeX,Scale=1}





\title{Quarto Manuals}


\shorttitle{Executable software operations manuals built from
independent Quarto pages}


\usepackage{etoolbox}









\authorsnames[{1},{1},{2,3},{4}]{Tinashe M.
Tapera,Chris,Danielle,Michelle}







\authorsaffiliations{
{Department of Psychology, Ana and Blanca's University},{Danielle's
Primary Affiliation},{Danielle's Secondary Affiliation},{Buffalo, NY }}




\leftheader{Tapera, Chris, Danielle and Michelle}



\abstract{to be fleshed out from bullet form below  \section{Impact
Statement} \noindent Scientific programming workflows are becoming
increasingly complex, while existing support mechanisms remain poorly
matched to that complexity. Traditional documentation is flexible but
non-executable, command-line wrappers are executable but often too
rigid, and current agentic systems are promising but not yet
sufficiently reliable for many high-stakes workflow setup tasks. We
propose Quarto Manuals, an executable documentation framework built on
Quarto that helps authors create flexible, scalable, and reproducible
workflow manuals composed of structured pages with explicit procedures
and checks. Powered by the quarto-emit backend, Quarto Manuals enable
authors to create executable documentation that allows users to
conveniently step through scientific computing workflows while
generating reproducible project artifacts and receiving immediate
feedback on whether each step has been implemented correctly. }

\keywords{APA Style, Quarto, Markdown, apaquarto}

\authornote{\par{\addORCIDlink{Tinashe M.
Tapera}{0000-0001-9080-5010}}\par{\addORCIDlink{Chris}{0000-0000-0000-0002}}\par{\addORCIDlink{Danielle}{0000-0000-0000-0003}}\par{\addORCIDlink{Michelle}{0000-0000-0000-0004}} 
\par{Carina Mengano is now at Generic University. }
\par{   The authors have no conflicts of interest to disclose.    Author
roles were classified using the Contributor Role Taxonomy (CRediT;
\href{https://credit.niso.org}{credit.niso.org}) as follows:  Tinashe M.
Tapera:   conceptualization, writing; Chris:   project
administration, formal analysis; Danielle:   formal
analysis, writing; Michelle:   writing, methodology, formal analysis}
\par{Correspondence concerning this article should be addressed
to Tinashe M. Tapera, Department of Psychology, Ana and Blanca's
University, 1234 Capital
St., Albany, NY 12084-1234, USA, Email: \href{mailto:ttapera@hsph.harvard.edu}{ttapera@hsph.harvard.edu}}
}

\makeatletter
\let\endoldlt\endlongtable
\def\endlongtable{
\hline
\endoldlt
}
\makeatother

\urlstyle{same}



\makeatletter
\@ifpackageloaded{caption}{}{\usepackage{caption}}
\AtBeginDocument{%
\ifdefined\contentsname
  \renewcommand*\contentsname{Table of contents}
\else
  \newcommand\contentsname{Table of contents}
\fi
\ifdefined\listfigurename
  \renewcommand*\listfigurename{List of Figures}
\else
  \newcommand\listfigurename{List of Figures}
\fi
\ifdefined\listtablename
  \renewcommand*\listtablename{List of Tables}
\else
  \newcommand\listtablename{List of Tables}
\fi
\ifdefined\figurename
  \renewcommand*\figurename{Figure}
\else
  \newcommand\figurename{Figure}
\fi
\ifdefined\tablename
  \renewcommand*\tablename{Table}
\else
  \newcommand\tablename{Table}
\fi
}
\@ifpackageloaded{float}{}{\usepackage{float}}
\floatstyle{ruled}
\@ifundefined{c@chapter}{\newfloat{codelisting}{h}{lop}}{\newfloat{codelisting}{h}{lop}[chapter]}
\floatname{codelisting}{Listing}
\newcommand*\listoflistings{\listof{codelisting}{List of Listings}}
\makeatother
\makeatletter
\makeatother
\makeatletter
\@ifpackageloaded{caption}{}{\usepackage{caption}}
\@ifpackageloaded{subcaption}{}{\usepackage{subcaption}}
\makeatother

% From https://tex.stackexchange.com/a/645996/211326
%%% apa7 doesn't want to add appendix section titles in the toc
%%% let's make it do it
\makeatletter
\xpatchcmd{\appendix}
  {\par}
  {\addcontentsline{toc}{section}{\@currentlabelname}\par}
  {}{}
\makeatother

%% Disable longtable counter
%% https://tex.stackexchange.com/a/248395/211326

\usepackage{etoolbox}

\makeatletter
\patchcmd{\LT@caption}
  {\bgroup}
  {\bgroup\global\LTpatch@captiontrue}
  {}{}
\patchcmd{\longtable}
  {\par}
  {\par\global\LTpatch@captionfalse}
  {}{}
\apptocmd{\endlongtable}
  {\ifLTpatch@caption\else\addtocounter{table}{-1}\fi}
  {}{}
\newif\ifLTpatch@caption
\makeatother

\begin{document}

\maketitle




\setlength\LTleft{0pt}



\ifdefined\Shaded\renewenvironment{Shaded}{\begin{snugshade}\begin{singlespace}}{\end{singlespace}\end{snugshade}}\fi


\ifdefined\cslhangindent\else\newlength{\cslhangindent}\fi
\ifdefined\csllabelwidth\else\newlength{\csllabelwidth}\fi
\setlength{\cslhangindent}{0.5in}


\subsection{Abstract}\label{abstract}

Software development is continuing to expand in both scale and
complexity, evidenced by the growth of global repository counts and
large-scale software
archives(\citeproc{ref-martinelliSoftwareHeritageActivity2026}{Martinelli,
2026};
\citeproc{ref-GithubInnovationGraph}{\textbf{GithubInnovationGraph?}})
At the same time, AI and agentic systems are reshaping how software is
produced, making it increasingly important to preserve accessibility,
transparency, and reproducibility in scientific software and data
science workflows.

Despite the central role of these workflows in academia and industry,
effective, reproducible documentation remains an overlooked and
under-served part of delivering high quality software
(\citeproc{ref-aghajaniSoftwareDocumentationIssues2019}{Aghajani et al.,
2019}). This issue is likely to become more acute as the volume and
sophistication of AI generated code submitted to public repositories
continues to increase
(\citeproc{ref-robbesAgenticMuchAdoption2026}{Robbes et al., 2026}).
Developers therefore need solutions beyond traditional documentation to
support workflow execution \_\_ address the limitations of traditional
software documentation to executing workflows in a reproducible and
verifiable manner as the landscape of software development continues to
evolve.

Conventional documentation can describe a workflow, but for complex,
multi-step, order-dependent tasks and projects, it is often insufficient
capturing and verifying critical runtime execution details such as
platform dependencies, configuration, versions, and relationships
between steps (\citeproc{ref-ebertGeneralConceptConsistent2015}{Ebert et
al., 2015}). Command-line interface (CLI) wrappers can improve
automation and reproducibility of a declarative workflow by providing a
directly executable interface for complex tasks
(\citeproc{ref-vanderaalstDeclarativeWorkflowsBalancing2009}{van der
Aalst et al., 2009}). However, developing and maintaining CLI wrappers
can be time-consuming and introduces a persistent trade-off between
control and flexibility: highly specialized wrappers can improve
consistency, but may be inflexible and difficult to maintain, while more
generic wrappers enable complexity by shifting the burden to the user
through large numbers of flags and configuration options
(\citeproc{ref-brackTenSimpleRules2022}{Brack et al., 2022};
\citeproc{ref-sadiqSpecificationValidationProcess2005}{Sadiq et al.,
2005}). Lastly, AI and agentic systems are promising, but their
reproducibility remains limited by non-determinism
(\citeproc{ref-siddiqLargeLanguageModels2025}{Siddiq et al., 2025}).
Furthermore, current practical agentic workflow implementations (such as
harnesses) still require careful human oversight to achieve acceptable
reliability (\citeproc{ref-agrawalCanAIConduct2026}{Agrawal et al.,
2026}).

Executable literate programming approaches can encode and enforce
complex workflow steps \emph{and} generate idempotent results. Quarto is
a modern, plain-text, multi-language, and multi-output literate
programming framework for creating reproducible documents and workflows
(\citeproc{ref-allaireQuarto2026}{Allaire et al., 2026}). Building on
Quarto with a carefully designed set of templates and extensions, we
introduce \textbf{Quarto Manuals}, a framework for executable software
manuals that guide users through a workflow as ordered, interactive
pages that combine explanation, code execution, and verification. Rather
than hiding workflow decisions and steps, Quarto Manuals keep them
visible and editable, making them ideal for research computing contexts
where local constraints and expert judgment often shape execution
parameters. The framework is powered by \texttt{quarto-emit}, a
lightweight backend extension that materializes conventional workflow
artifacts from manual pages when needed. We demonstrate the approach
with three examples of increasing sophistication: a simple manual, an
intermediate manual, and a complex manual.

For authors, Quarto Manuals provide a structured yet adaptable way to
turn recurring workflows into reusable manuals; for operators, they
provide a stepwise, testable process for generating reproducible setup
artifacts and confirming progress throughout execution in a structured
and auditable manner.

\subsection{Feature Roadmap}\label{feature-roadmap}

\begin{itemize}
\tightlist
\item[$\boxtimes$]
  abstract as prose
\item[$\square$]
  stylize each block so it is visually distinct
\item[$\square$]
  use Lua to print out useful render messages when blocks are being
  processed
\item[$\square$]
  negative reinforcement: identify and correct misuse of the manual
  visually
\item[$\square$]
  setup a testing suite in a separate repo
\item[$\square$]
  can Lua be used to connect checks and prereqs across pages? Maybe?
\end{itemize}

\section{References}\label{references}

References marked with an asterisk indicate studies included in the
meta-analysis.

\protect\phantomsection\label{refs}
\begin{CSLReferences}{1}{0}
\bibitem[\citeproctext]{ref-aghajaniSoftwareDocumentationIssues2019}
Aghajani, E., Nagy, C., Vega-Márquez, O. L., Linares-Vásquez, M.,
Moreno, L., Bavota, G., \& Lanza, M. (2019). Software {Documentation
Issues Unveiled}. \emph{2019 {IEEE}/{ACM} 41st {International
Conference} on {Software Engineering} ({ICSE})}, 1199--1210.
\url{https://doi.org/10.1109/ICSE.2019.00122}

\bibitem[\citeproctext]{ref-agrawalCanAIConduct2026}
Agrawal, S., Anadkat, H. B., Athimoolam, K. K., Bhardwaj, H., Chowdhury,
T., Gao, S., Kamat, P. K., Makwana, V., Shariff, M. H., Badkul, A., Xie,
L., \& Sinitskiy, A. V. (2026, January 6). \emph{Can {AI Conduct
Autonomous Scientific Research}? {Case Studies} on {Two Real-World
Tasks}}. \url{https://doi.org/10.64898/2026.01.05.697809}
(Pre-published)

\bibitem[\citeproctext]{ref-allaireQuarto2026}
Allaire, J. J., Teague, C., Scheidegger, C., Xie, Y., Dervieux, C., \&
Woodhull, G. (2026). \emph{Quarto} (Version 1.10) {[}Computer
software{]}. \url{https://doi.org/10.5281/zenodo.5960048}

\bibitem[\citeproctext]{ref-brackTenSimpleRules2022}
Brack, P., Crowther, P., Soiland-Reyes, S., Owen, S., Lowe, D.,
Williams, A. R., Groom, Q., Dillen, M., Coppens, F., Grüning, B.,
Eguinoa, I., Ewels, P., \& Goble, C. (2022). Ten simple rules for making
a software tool workflow-ready. \emph{PLOS Computational Biology},
\emph{18}(3), e1009823.
\url{https://doi.org/10.1371/journal.pcbi.1009823}

\bibitem[\citeproctext]{ref-ebertGeneralConceptConsistent2015}
Ebert, P., Müller, F., Nordström, K., Lengauer, T., \& Schulz, M. H.
(2015). A general concept for consistent documentation of computational
analyses. \emph{Database: The Journal of Biological Databases and
Curation}, \emph{2015}, bav050.
\url{https://doi.org/10.1093/database/bav050}

\bibitem[\citeproctext]{ref-martinelliSoftwareHeritageActivity2026}
Martinelli, N. (2026, January 16). \emph{Software {Heritage Activity
Report}: 2025}. Software Heritage.
\url{https://www.softwareheritage.org/2026/01/16/software-heritage-activity-report-2025/}

\bibitem[\citeproctext]{ref-robbesAgenticMuchAdoption2026}
Robbes, R., Matricon, T., Degueule, T., Hora, A., \& Zacchiroli, S.
(2026). Agentic {Much}? {Adoption} of {Coding Agents} on {GitHub}.
\emph{ACM Transactions on Software Engineering and Methodology}.
\url{https://doi.org/10.1145/3822180}

\bibitem[\citeproctext]{ref-sadiqSpecificationValidationProcess2005}
Sadiq, S. W., Orlowska, M. E., \& Sadiq, W. (2005). Specification and
validation of process constraints for flexible workflows.
\emph{Information Systems}, \emph{30}(5), 349--378.
\url{https://doi.org/10.1016/j.is.2004.05.002}

\bibitem[\citeproctext]{ref-siddiqLargeLanguageModels2025}
Siddiq, M. L., Islam-Gomes, A., Sekerak, N., \& Santos, J. C. S. (2025,
November 29). \emph{Large {Language Models} for {Software Engineering}:
{A Reproducibility Crisis}}.
\url{https://doi.org/10.48550/arXiv.2512.00651}

\bibitem[\citeproctext]{ref-vanderaalstDeclarativeWorkflowsBalancing2009}
van der Aalst, W. M. P., Pesic, M., \& Schonenberg, H. (2009).
Declarative workflows: {Balancing} between flexibility and support.
\emph{Computer Science - Research and Development}, \emph{23}(2),
99--113. \url{https://doi.org/10.1007/s00450-009-0057-9}

\end{CSLReferences}






\end{document}
